\section{Priorisierung und Abgrenzung}
Die priorisierte Bewertung der Anforderungen ermöglicht eine strukturierte Umsetzung.
Dabei werden zentrale Funktionen gegenüber optionalen Erweiterungen abgegrenzt.

\begin{table}[H]
\centering
\begin{tabular}{p{3cm} p{3cm} p{6cm}}
\textbf{ID} & \textbf{Priorität} & \textbf{Begründung} \\
\hline
F1–F3 & Hoch & Kernfunktion der Wortuhr \\
F4–F5 & Hoch & Erweiterung mit hohem Mehrwert \\
F6–F9 & Mittel & Zusatzfunktionen zur Informationsdarstellung \\
F10–F12 & Hoch & Notwendig für Vernetzung und Bedienung \\
N1–N3 & Hoch & Grundlegende Systemeigenschaften \\
N4–N6 & Mittel & Qualitäts- und Komfortmerkmale \\
\end{tabular}
\caption{Priorisierung der Anforderungen}
\end{table}

Die Abgrenzung erfolgt insbesondere hinsichtlich des Entwicklungsumfangs.
Ziel ist die Realisierung eines funktionsfähigen Prototyps.
Aspekte wie industrielle Serienfertigung, normgerechte Zertifizierung oder wirtschaftliche Optimierung werden nicht betrachtet.

Die priorisierte Struktur dient als Grundlage für die anschließende Systemkonzeption und Implementierung.
% %%%%%%%%%%%%Old Version
% Die Anforderungen F1 und F2 stellen Muss-Kriterien dar. Ohne diese Funktionen erfüllt das System nicht die grundlegende Zielsetzung einer Wortuhr.
% Die Anforderungen F4 und F5 gelten als Soll-Kriterien. Sie erweitern den Funktionsumfang, sind jedoch nicht zwingend für die Kernfunktion erforderlich.
% Eine Sprachsteuerung, eine Integration in externe Smart-Home-Plattformen sowie eine mobile Applikation werden im Rahmen dieser Arbeit nicht umgesetzt. Diese Funktionen erfordern zusätzlichen Entwicklungsaufwand und liegen außerhalb des definierten Projektumfangs.
% Die Priorisierung dient der strukturierten Umsetzung innerhalb des vorgegebenen Zeitrahmens. Sie ermöglicht eine fokussierte Entwicklung der zentralen Funktionen und schafft eine Grundlage für mögliche zukünftige Erweiterungen.
